\subsection{Two perturbative auxiliary results}\label{app:auxiliary-results}

In this appendix we prove two perturbative results used in the proof of
Theorem~\ref{thm:local-perturbative}. The first is a free-boundary stability
statement around a strictly negative obstacle-type solution, in the spirit of
the perturbative one-phase results of Isakov~\cite[Chapter~5]{Isakov} and the
classical free-boundary literature on quadrature domains
\cite{Aharonov-Schiffer-Zalcman,Karp,Gustafsson,Shahgholian,Gustafsson-Shahgholian}.
The second is a local realization theorem for weighted quadrature domains near
a reference domain. Both statements are tailored to the setting of the present
paper.

\subsubsection{A perturbative free-boundary theorem}\label{app:isakov-type}

We begin with a perturbative result for the one-phase problem
\[
-\Delta u = f\,\mathbf{1}_{\Omega},
\qquad
u=0 \ \text{on } \C\setminus \Omega,
\]
near a reference pair \((\Omega_0,u_0)\). The reader interested in more general
statements of this type may consult Chapter~5 of~\cite{Isakov}.

Let \(\Omega_0\subset \C\simeq \R^2\) be a bounded domain of class \(C^2\). For
\(\rho>0\) sufficiently small, every \(C^2\) domain \(\Omega\) which is
\(\rho\)-close to \(\Omega_0\) (in the Hausdorff sense, or equivalently in the
sense of \(C^2\)-graphs over \(\partial\Omega_0\)) may be written uniquely as
a normal graph over \(\partial\Omega_0\):
\[
\partial\Omega
=
\{x+\varphi(x)\nu(x):x\in \partial\Omega_0\},
\]
where \(\nu\) denotes the outer unit normal to \(\partial\Omega_0\), and
\(\varphi\in C^2(\partial\Omega_0)\) has small norm. We shall write
\(\Omega=\Omega_\varphi\). Throughout this subsection we keep a H\"older
exponent \(\alpha\in(0,1)\) fixed; the choice of \(C^5\) regularity on the
right-hand side \(f\) is dictated by the requirement that the implicit
function theorem close in the space of \(C^2\) (equivalently
\(C^{2,\alpha}\)) boundaries: \(C^5\) data on \(f\) and \(C^2\) regularity of
\(\partial\Omega_0\) ensure that the solution operator
\((\varphi,f)\mapsto u_{\varphi,f}\circ\Phi_\varphi\) is of class \(C^1\) into
\(C^{2,\alpha}(\overline{\Omega_0})\) on a small neighbourhood of
\((0,f_0)\).

\begin{proposition}[\cite{Isakov}] \label{prop:isakov-tailored}
Let \(\Omega_0\subset \C\) be a bounded \(C^2\) domain, and let
\(f_0\in C^5(\overline{\Omega'})\) for some open neighborhood
\(\Omega'\supset \overline{\Omega_0}\), with
\[
f_0>0 \qquad \text{on } \partial\Omega_0.
\]
Assume there exists \(u_0\in C^{2,\alpha}(\overline{\Omega_0})\) for some
\(\alpha\in(0,1)\) such that
\begin{equation}\label{eq:appendix-reference-problem}
\left\{
\begin{aligned}
-\Delta u_0 &= f_0 &&\text{in }\Omega_0,\\
u_0&=0 &&\text{on }\partial\Omega_0,\\
u_0&<0 &&\text{in }\Omega_0.
\end{aligned}
\right.
\end{equation}
Then there exist \(\varepsilon>0\) and \(\rho>0\) such that the following holds.

Whenever \(f\in C^5(\overline{\Omega'})\) satisfies
\[
\|f-f_0\|_{C^5(\overline{\Omega'})}<\varepsilon,
\]
there exists a unique \(C^2\) function
\(\varphi\in C^2(\partial\Omega_0)\), with \(\|\varphi\|_{C^2}<\rho\), such
that for \(\Omega_f:=\Omega_\varphi\) the Dirichlet problem
\begin{equation}\label{eq:appendix-perturbed-problem}
\left\{
\begin{aligned}
-\Delta u_f &= f &&\text{in }\Omega_f,\\
u_f&=0 &&\text{on }\partial\Omega_f
\end{aligned}
\right.
\end{equation}
has a solution \(u_f\in C^{2,\alpha}(\overline{\Omega_f})\) satisfying
\[
u_f<0 \qquad \text{in }\Omega_f.
\]
Moreover, the map
\[
f\mapsto \varphi(f)
\]
is \(C^1\) from a neighborhood of \(f_0\) in \(C^5(\overline{\Omega'})\) into
\(C^2(\partial\Omega_0)\); in particular, \(\varphi\) depends continuously on
\(f\) in the indicated norms.
\end{proposition}

\begin{proof}
We divide the argument into several steps; the structure follows that of
Corollary~5.1.4 of~\cite{Isakov}, adapted to the present formulation.

\medskip
\noindent\emph{Step 1: nondegeneracy of the reference solution.}
By assumption, \(f_0>0\) on the compact set \(\partial\Omega_0\). Since
\(f_0\in C^5(\overline{\Omega'})\) is, in particular, continuous on
\(\overline{\Omega_0}\), there exist constants \(c_0>0\) and \(\eta>0\) such
that
\[
f_0(x)\ge c_0
\qquad \text{for every } x\in \overline{\Omega_0}
\text{ with }\operatorname{dist}(x,\partial\Omega_0)<\eta.
\]
Combined with \(u_0<0\) in \(\Omega_0\), \(u_0=0\) on \(\partial\Omega_0\),
and \(-\Delta u_0=f_0\ge c_0>0\) in this neighborhood of \(\partial\Omega_0\),
the Hopf boundary point lemma (in the standard application discussed in the
perturbative one-phase setting of~\cite[Chapter~5]{Isakov}) yields
\[
\partial_\nu u_0(x)>0
\qquad \text{for every }x\in \partial\Omega_0,
\]
where \(\nu\) is the outer unit normal. Since \(\partial\Omega_0\) is compact
and \(u_0\in C^{2,\alpha}(\overline{\Omega_0})\), the function
\(x\mapsto \partial_\nu u_0(x)\) is continuous on \(\partial\Omega_0\), and we
obtain a uniform lower bound
\begin{equation}\label{eq:appendix-hopf-lower-bound}
\partial_\nu u_0\ge \kappa_0
\qquad\text{on }\partial\Omega_0,
\end{equation}
for some constant \(\kappa_0>0\).

\medskip
\noindent\emph{Step 2: reduction to a boundary equation.}
Fix a tubular neighborhood
\[
\mathcal N:=\{x+t\nu(x):x\in \partial\Omega_0,\ |t|<\delta\}
\]
on which the normal coordinates are well-defined and the map
\((x,t)\mapsto x+t\nu(x)\) is a \(C^2\) diffeomorphism onto its image. For
\(\varphi\in C^2(\partial\Omega_0)\) sufficiently small in \(C^2\), the
perturbed boundary
\[
\Gamma_\varphi:=\{x+\varphi(x)\nu(x):x\in \partial\Omega_0\}
\]
bounds a \(C^2\) domain \(\Omega_\varphi\subset \Omega'\), and the normal
graph map
\[
\Phi_\varphi:\overline{\Omega_0}\longrightarrow \overline{\Omega_\varphi},
\]
constructed by extending \(\varphi\) to a \(C^2\) function
\(\widetilde\varphi\) on a neighbourhood of \(\partial\Omega_0\) and setting
\(\Phi_\varphi(y)=y+\widetilde\varphi(y)\nu(y)\) on
\(\mathcal N\cap \Omega_0\) and the identity outside, is a \(C^2\)
diffeomorphism with \(\Phi_0=\mathrm{Id}\), depending in a \(C^1\) fashion on
\(\varphi\) in the \(C^2\)-topology.

For each pair \((\varphi,f)\) with \(\varphi\) small in
\(C^2(\partial\Omega_0)\) and \(f\) close to \(f_0\) in
\(C^5(\overline{\Omega'})\), let \(u_{\varphi,f}\) denote the unique solution
of
\[
\left\{
\begin{aligned}
-\Delta u_{\varphi,f} &= f &&\text{in }\Omega_\varphi,\\
u_{\varphi,f}&=0 &&\text{on }\Gamma_\varphi.
\end{aligned}
\right.
\]
Pulling back via \(\Phi_\varphi\), the function
\(v_{\varphi,f}:=u_{\varphi,f}\circ \Phi_\varphi\) satisfies a transformed
boundary-value problem of the form
\[
\left\{
\begin{aligned}
L_\varphi v_{\varphi,f}&=f\circ \Phi_\varphi &&\text{in }\Omega_0,\\
v_{\varphi,f}&=0 &&\text{on }\partial\Omega_0,
\end{aligned}
\right.
\]
where \(L_\varphi\) is a second-order uniformly elliptic operator with
\(C^{0,\alpha}\) (in fact \(C^1\)) coefficients depending in a \(C^1\)
fashion on \(\varphi\), and with \(L_0=-\Delta\). Standard Schauder theory
for elliptic boundary value problems (see e.g.~\cite[Ch.~5]{Isakov}) then
implies that the map
\[
(\varphi,f)\mapsto v_{\varphi,f}=u_{\varphi,f}\circ \Phi_\varphi
\]
is of class \(C^1\) from a neighborhood of \((0,f_0)\) in
\(C^2(\partial\Omega_0)\times C^5(\overline{\Omega'})\) into
\(C^{2,\alpha}(\overline{\Omega_0})\); the choice of \(C^5\) data on the
right-hand side guarantees the requisite regularity of the coefficients of
\(L_\varphi\) and of the source term \(f\circ\Phi_\varphi\) in this Schauder
framework, and in particular that the linearised problems are well-posed in
\(C^{2,\alpha}\).

Now define the nonlinear boundary operator
\[
\mathcal F(\varphi,f)(x)
:=
u_{\varphi,f}\bigl(x+\varphi(x)\nu(x)\bigr),
\qquad x\in \partial\Omega_0;
\]
here we regard \(u_{\varphi,f}\) as a \(C^{2,\alpha}\) function defined in a
neighbourhood of \(\overline{\Omega_\varphi}\) (which contains a fixed
neighbourhood of \(\overline{\Omega_0}\) for \(\varphi\) small in \(C^2\))
and evaluate it at the point \(x+\varphi(x)\nu(x)\), which lies on
\(\Gamma_\varphi\). Since \(u_{\varphi,f}\) vanishes on \(\Gamma_\varphi\), the
condition that \(\Gamma_\varphi\) be the free boundary of
\eqref{eq:appendix-perturbed-problem} is precisely
\begin{equation}\label{eq:appendix-free-boundary-equation}
\mathcal F(\varphi,f)=0.
\end{equation}
By the regularity statement above, \(\mathcal F\) is \(C^1\) from a
neighborhood of \((0,f_0)\) into \(C^{2,\alpha}(\partial\Omega_0)\).

\medskip
\noindent\emph{Step 3: computation of the linearization in \(\varphi\).}
For the purpose of the implicit function theorem we replace \(\mathcal F\)
by the equivalent boundary functional
\[
\widetilde{\mathcal F}(\varphi,f)(x)
:=
u_{0,f}\bigl(x+\varphi(x)\nu(x)\bigr),
\qquad x\in\partial\Omega_0,
\]
where \(u_{0,f}\) denotes the solution of \(-\Delta u=f\) on the
\emph{fixed} reference domain \(\Omega_0\), with \(u_{0,f}=0\) on
\(\partial\Omega_0\), extended to a fixed neighbourhood of
\(\overline{\Omega_0}\) by a bounded linear extension operator
\(E\colon C^{2,\alpha}(\overline{\Omega_0})\to C^{2,\alpha}(\R^2)\) (such
an extension exists by standard arguments, e.g.\ Stein-type extension or
local reflection across the \(C^2\) boundary \(\partial\Omega_0\)). For
\((\varphi,f)\) close to \((0,f_0)\) the equations
\[
\widetilde{\mathcal F}(\varphi,f)=0
\qquad\text{and}\qquad
\bigl(\Gamma_\varphi\text{ is the free boundary of
}\eqref{eq:appendix-perturbed-problem}\bigr)
\]
cut out the same set of \((\varphi,f)\): both express that
\(\Gamma_\varphi\) is the zero level set of the solution corresponding to
\(f\). The advantage of \(\widetilde{\mathcal F}\) is that the underlying
solution \(u_{0,f}\) does not depend on \(\varphi\), so its derivative in
\(\varphi\) reduces to a pure shape derivative of a fixed function. From
now on we abuse notation and continue to write \(\mathcal F\) for
\(\widetilde{\mathcal F}\).

Let \(\psi\in C^2(\partial\Omega_0)\), extend \(\psi\) to a \(C^2\) function
in a neighbourhood of \(\partial\Omega_0\), and consider the family
\(\varphi_\varepsilon=\varepsilon\psi\), \(\varepsilon\in\R\). Using
\(u_0(x)=0\) for \(x\in\partial\Omega_0\) and Taylor-expanding \(u_0\) at
first order in the normal direction,
\[
\mathcal F(\varepsilon\psi,f_0)(x)
=
u_0\bigl(x+\varepsilon\psi(x)\nu(x)\bigr)
=
u_0(x)+\varepsilon\,\psi(x)\,\partial_\nu u_0(x)+o(\varepsilon)
=
\varepsilon\,\psi(x)\,\partial_\nu u_0(x)+o(\varepsilon),
\]
uniformly in \(x\in \partial\Omega_0\). Hence
\[
D_\varphi \mathcal F(0,f_0)[\psi]
=
\psi\,\partial_\nu u_0
\qquad \text{on }\partial\Omega_0.
\]
This is the standard Hadamard formula in the present Dirichlet setting:
moving the boundary in the normal direction by \(\psi\) changes the boundary
trace of the (fixed) solution at first order by \(\psi\,\partial_\nu u_0\).

In view of \eqref{eq:appendix-hopf-lower-bound}, the multiplication operator
\[
\psi\mapsto \psi\,\partial_\nu u_0
\]
is invertible, with continuous inverse \(g\mapsto g/\partial_\nu u_0\). Since
\(\partial_\nu u_0\in C^{1,\alpha}(\partial\Omega_0)\) (being the boundary
trace of a \(C^{1,\alpha}\) vector field on a \(C^2\) curve, derived from
the \(C^{2,\alpha}\) function \(u_0\)) and is bounded below by
\(\kappa_0>0\), this operator is an isomorphism of \(C^2(\partial\Omega_0)\)
onto itself, and likewise of \(C^{2,\alpha}(\partial\Omega_0)\) onto itself.

\medskip
\noindent\emph{Step 4: application of the implicit function theorem.}
We apply the implicit function theorem in Banach spaces to the \(C^1\) map
\[
\mathcal F:U\subset C^2(\partial\Omega_0)\times C^5(\overline{\Omega'})
\longrightarrow C^{2,\alpha}(\partial\Omega_0)
\]
at the point \((0,f_0)\). We have \(\mathcal F(0,f_0)=0\), and by Step~3 the
partial derivative \(D_\varphi\mathcal F(0,f_0)\) is an isomorphism onto its
image. The implicit function theorem therefore yields \(\varepsilon>0\),
\(\rho>0\), and a unique \(C^1\) map
\[
f\mapsto \varphi(f)\in C^2(\partial\Omega_0),
\qquad
\|\varphi(f)\|_{C^2}<\rho,
\]
defined for \(\|f-f_0\|_{C^5}<\varepsilon\), such that
\[
\mathcal F(\varphi(f),f)=0,
\qquad \varphi(f_0)=0.
\]
Equivalently, the corresponding domain
\(\Omega_f:=\Omega_{\varphi(f)}\) has the property that the solution
\(u_f:=u_{\varphi(f),f}\) of \eqref{eq:appendix-perturbed-problem} vanishes
on \(\partial\Omega_f\). This already gives existence and uniqueness of
\(\varphi(f)\) inside the small \(C^2\)-ball of radius \(\rho\) about \(0\).

\medskip
\noindent\emph{Step 5: preservation of the sign.}
It remains to show that, possibly after shrinking \(\varepsilon\) and
\(\rho\), the solution \(u_f\) is strictly negative on the whole of
\(\Omega_f\).

By Step~2, the map \((\varphi,f)\mapsto u_{\varphi,f}\circ \Phi_\varphi\) is
\(C^1\) (in particular continuous) from a neighbourhood of \((0,f_0)\) into
\(C^{2,\alpha}(\overline{\Omega_0})\). By the embedding
\(C^{2,\alpha}\hookrightarrow C^1\),
\begin{equation}\label{eq:appendix-c1-convergence}
\bigl\|u_f\circ \Phi_{\varphi(f)}-u_0\bigr\|_{C^1(\overline{\Omega_0})}
\longrightarrow 0
\qquad \text{as } f\to f_0
\text{ in } C^5(\overline{\Omega'}).
\end{equation}
To deduce strict negativity from this \(C^1\)-closeness, we split
\(\overline{\Omega_0}\) into a fixed interior region and a thin collar
around the boundary, where additional care is needed because \(u_0=0\) on
\(\partial\Omega_0\).

Fix \(\delta_0>0\) so small that the inner collar
\[
\mathcal N_{\delta_0}:=\{x-t\nu(x):x\in \partial\Omega_0,\ 0<t<\delta_0\}
\]
is contained in \(\Omega_0\) and the map \((x,t)\mapsto x-t\nu(x)\) is a
diffeomorphism of \(\partial\Omega_0\times (0,\delta_0)\) onto
\(\mathcal N_{\delta_0}\). By \eqref{eq:appendix-hopf-lower-bound} and the
\(C^1\)-regularity of \(u_0\) up to the boundary, after possibly shrinking
\(\delta_0\) we obtain the quantitative collar estimate
\[
u_0(x-t\nu(x))\le -\frac{\kappa_0}{2}\,t
\qquad
\text{for }x\in \partial\Omega_0,\ 0<t<\delta_0.
\]
Set
\[
K_{\delta_0}
:=
\bigl\{y\in \Omega_0:\operatorname{dist}(y,\partial\Omega_0)\ge \delta_0/2\bigr\},
\]
which is a compact subset of \(\Omega_0\) on which, by continuity and strict
negativity of \(u_0\) in \(\Omega_0\), one has \(u_0\le -m_0\) for some
\(m_0>0\).

By \eqref{eq:appendix-c1-convergence}, for \(f\) sufficiently close to
\(f_0\) in \(C^5(\overline{\Omega'})\) the function
\(v_f:=u_f\circ\Phi_{\varphi(f)}\) satisfies
\[
\|v_f-u_0\|_{C^1(\overline{\Omega_0})}
<\min\!\left\{\frac{\kappa_0}{4},\frac{m_0}{2}\right\}.
\]
On \(K_{\delta_0}\) this immediately gives
\(v_f\le -m_0+m_0/2=-m_0/2<0\). On the closed collar
\(\overline{\mathcal N_{\delta_0}}\), the \(C^1\)-closeness combined with
\eqref{eq:appendix-hopf-lower-bound} gives, on \(\partial\Omega_0\),
\[
\partial_\nu v_f
\ge
\partial_\nu u_0-\|v_f-u_0\|_{C^1}
\ge \kappa_0-\frac{\kappa_0}{4}=\frac{3\kappa_0}{4}.
\]
By uniform continuity of \(\nabla v_f\) and \(\nabla u_0\) on
\(\overline{\Omega_0}\) — recall both lie in a fixed bounded set of
\(C^{1,\alpha}(\overline{\Omega_0})\) by the Schauder estimate, hence
in a precompact subset of \(C^1(\overline{\Omega_0})\) — after a further
shrinking of \(\delta_0\) (independent of \(f\)) we have, for every
\(x\in\partial\Omega_0\) and \(0<t<\delta_0\),
\[
\partial_\nu v_f(x-t\nu(x))\ge \frac{\kappa_0}{2}.
\]
Since \(v_f=0\) on \(\partial\Omega_0\), the fundamental theorem of calculus
applied along inward normal segments yields
\[
v_f(x-t\nu(x))
=
v_f(x)-\int_0^{t} \partial_\nu v_f(x-s\nu(x))\,ds
\le -\frac{\kappa_0}{2}\,t<0
\qquad \text{for } x\in \partial\Omega_0,\ 0<t<\delta_0.
\]
Hence \(v_f<0\) on \(\mathcal N_{\delta_0}\) as well, and combining with the
estimate on \(K_{\delta_0}\) we conclude that
\(u_f\circ\Phi_{\varphi(f)}<0\) throughout \(\Omega_0\). Since
\(\Phi_{\varphi(f)}\) is a diffeomorphism from \(\Omega_0\) onto
\(\Omega_f\), this is equivalent to \(u_f<0\) throughout \(\Omega_f\).

This proves existence. Uniqueness among sufficiently small \(C^2\) normal
graphs follows from the uniqueness statement in the implicit function theorem
applied in Step~4.
\end{proof}

